\section{Esoteric Scholasticism}

\begin{quotex}
The first man was of the earth, earthly: the second man, from heaven, heavenly. Such as is the earthly, such also are the earthly: and such as is the heavenly, such also are they that are heavenly. Therefore as we have borne the image of the earthly, let us bear also the image of the heavenly. \flright{\textsc{1 Cor 15:47-49}}

Scholasticism is in the highest degree a product of human ingenuity. In it the logical faculty celebrated its greatest triumphs. One who aims to elaborate concepts in their sharpest and clearest contours should serve an apprenticeship with the Scholastics. It is they who provide the highest schooling for the technique of thinking. They have an incomparable agility in moving in the field of pure thought. \flright{\textsc{Rudolf Steiner}, \emph{Eleven European Mystics} (GA7)}

Dear Unknown Friend, do not scorn mediaeval scholasticism. It is, in truth, as beautiful, as venerable and as inspiring as the great cathedrals that we have inherited from the Middle Ages. To it we owe a number of masterpieces of thought — thought in the light of faith. And, like all true masterpieces, those of mediaeval scholasticism are beneficial. They heal the disorientated, feverous and confused soul. \flright{\textsc{Valentin Tomberg}, \emph{Meditations on the Tarot}}

\end{quotex}
Scholasticism deals with the most important concepts such as God, the soul, freedom, immortality, salvation, good and evil. The foundation for thinking begins with the seven liberal arts: grammar, logic, rhetoric, arithmetic, geometry, music, astronomy. Compare that to the stunted interests of philosophers nowadays.

Although knowledge begins with the senses, it is incomplete until the idea is known. Otto Willmann described the Scholastic process to higher understanding:

\begin{quotex}
Our spirit, associated with the body as it is in earthly life, is primarily directed toward the surrounding world of matter, but \textbf{focused upon the spiritual in it; that is, the essences, natures, and forms of things}, the elements of existence which are akin to it and provide it with the rungs by which it ascends to the supra-sensory. The field of our cognition is thus the realm of experience, but we should learn to understand what it offers, penetrate to its sense and idea, and thereby open to ourselves the world of ideas. 

\end{quotex}
Ultimately the goal is self-knowledge, not merely the image of his self. Karl Werner warns about being satisfied with the latter rather than the former.

\begin{quotex}
In time man has no perception of his self, the hidden foundation of his spiritual nature and life; … he will never be able to look at himself; for either, forever estranged from God, he will find in himself only a bottomless dark abyss and endless emptiness, or he will, blessed in God, and turning his gaze inward, find only God, Whose sun of grace shines within him, and Whose \emph{image} reflects itself in the spiritual traits of his nature. 

\end{quotex}
But \textbf{St Augustine} understood that the actual God image in man is Christ rather than that self-image. The God image is in the \emph{anima rationalis}, the intellectual soul, which distinguishes the human from the beast; in effect, then, it is supernatural, not natural. He elaborates:

\begin{quotex}
The God-image is within, not in the body. … Where the understanding is, where the mind is, where the power of investigating truth is, there God has his image. … where man knows himself to be made after the image of God, there he knows there is something more in him than is given to the beasts. 

\end{quotex}
That sums up the Scholastic project. Start with the understanding and work up to the concept of God. That is the path of the intellect. Then follow up with the Mystical Impulse.

\paragraph{The Mystical Impulse}
I cannot express this better than the Unknown Friend:

\begin{quotex}
No more is it true that the mystical impulse from the end of the thirteenth and into the seventeenth century was purely and simply a reaction against the ``dry intellectualism" of scholasticism. No, \textbf{the flowering of mysticism during this epoch was the fruit and the result of scholasticism}.

St. Thomas towards the end of his life arrived at mystical contemplation of God and the spiritual world and said, on returning from this ecstasy, that his written works now appeared to him ``like straw". Indeed, he wrote nothing after this. The believing thinker thus became a seeing mystic. And this transformation did not take place in spite of his work of scholastic thought, but rather thanks to it—as its fruit and its crowning glory.

Now, what happened to St. Thomas Aquinas also happened to a group of individuals who formed the crest of the wave of scholasticism. Just as St. Thomas, through scholastic reasoning, arrived at contemplation, so did part of advanced scholasticism arrive at mysticism, i.e. at the \emph{aim} of scholasticism, which is \emph{intuition} or \textbf{the state of union of faith and intelligence}. 

\end{quotex}
\paragraph{Renaissance Mystics}
Valentin Tomberg mentions three mystics in particular:

\begin{quotex}
Meister Eckhart, Ruysbroeck, or, lastly, St. John of the Cross are spirits amongst whom you will search in vain for a spirit of opposition to scholasticism. For them also it was true that scholasticism was ``like straw", but they knew at the same time from their own experience that this straw proved to be an excellent combustible. They certainly surpassed scholasticism, but they did so by attaining its aim. \textbf{For the aim of scholastic endeavour is contemplation, and it is mysticism which is the fruit of the scholastic tree}. 

\end{quotex}
Rudolf Steiner has a similar point of view:

\begin{quotex}
One can therefore interpret the activity of the mystical theologians Eckhart, Tauler, Suso and their companions by saying: They were inspired by the content of the teachings of the Church, which is contained in theology, but had been reinterpreted, to bring forth a similar content out of themselves anew as an \textbf{inner experience}. 

\end{quotex}
The following list combines suggestions from Tomberg and Steiner. The recovery of esoteric or mystical Scholasticism is the antidote to the scientific mechanistic dualism that passes for knowledge today.

\begin{itemize}
\item Meister Eckhart 
\item Johannes Tauler 
\item Heinrich Suso 
\item Jan van Ruysbroeck 
\item Nicolas of Cusa 
\item Cornelius Agrippa 
\item Paracelsus 
\item Valentin Weigel 
\item Jacob Boehme 
\item Giordano Bruno 
\item Angelus Silesius 
\item St John of the Cross 
\end{itemize}


\flrightit{Posted on 2020-07-14 by Cologero }

\begin{center}* * *\end{center}

\begin{footnotesize}\begin{sffamily}



\texttt{Adil on 2020-07-15 at 10:55 said: }

And there goes the myth of dry scholastic intellectualism, that ``mysteriously" and supposedly paved ground for atheism (according to some commentators).


\hfill

\texttt{Michael M on 2020-07-16 at 09:41 said: }

@Adil

Perhaps the opening of mysticism is just ``too much" to recognize for moderns considering the scope of the effort to produce such an effect in the World. Incredibly hard to verify, we specifically lack this higher quality tied to overcoming the discursive intellect. What is typically the argument those commentators make? Are they just seeing their own image in the Scholastics? Maybe given that Man now doesn't see his true image as a reflection of God but of the highest truth as man, he can only look back and see himself in all things.


\hfill

\texttt{Adil on 2020-07-16 at 12:09 said: }

Michael M

I've heard the argument circulate in some Eastern Orthodox corners of the net. They seem to push forward that the `backwards' empirical reasoning of Aquinas from sensory input to God would leave this First Mover bereft of any positive characteristics (circumventing the mystical), thus begging the question for the rise of atheism. The opposite conclusion of this article, which claims scholasticism laid the groundwork for mystical integration. Which was indeed confirmed by Aquinas's own breakthrough.

``Maybe given that Man now doesn't see his true image as a reflection of God but of the highest truth as man, he can only look back and see himself in all things."

It seems this is a corruption of the Christ impulse, in which spirit seeks to realize and transform itself through man – but backfires in solipsism and egotism. The only mystical breakthrough in such a state is that of the temporary intoxication.


\hfill

\texttt{Cologero on 2020-07-16 at 17:09 said: }

I don't know about those corners.

Empirical reasoning from phenomena would look more like a design argument, which Thomas rejects. From the senses, you then develop and understanding of essences. Not at all ``empirical". And there is indeed some positive content.

If God is not the Unmoved Mover, then that means he is not changeless.

If God is not the First Cause, then He must have had some help to create the universe.

That's an odd argument from those corners since they deny the possibility of the knowledge of God's essence. Actually, Thomas is not all that distant from Palamas, since he also denies that possibility (at least in this life). Hence, we can know God only through his ``activities", which is the Latin equivalent of the Greek-derived ``energies". (the root Act in Latin translates the Greek erg). If those are not empirical, what exactly are they?


\hfill

\texttt{Adil on 2020-07-18 at 04:44 said: }

I believe the problem such an Orthodox critique posits is precisely that there is no direct experience of the divine in this world. That is, since according to Thomistic divine simplicity, God's nature cannot be divided, the essences turn into mere `phantoms' of God's mind, which is ultimately unknowable. Therefore, there is a fundamental gap between God and this world, which can't be directly penetrated by the uncreated divine grace. This, supposedly, would be fertile breeding ground for atheistic ideas as well as de-sacralizing nature and leading further into Protestant dualism. More generally, I believe it's a critique of Aristotle's ``materialism" (denial of the pure realm of forms).


\hfill

\texttt{Cologero on 2020-07-21 at 07:33 said: }

Divine simplicity is also accepted by the Orthodox. Aristotle is not a materialist. He emphasizes the immanent aspect of ``forms", while Plato emphasizes the transcendent aspect. There is no ``realm of forms" unless you have a better geography book than mine; the forms are in the Mind of God. If there were a ``direct experience of the divine in this world", then scientists would have discovered it.

The conclusion is, as we point out in this essay, that Thomas provides a ``fertile breeding ground" for a direct, mystical intuition of God. However, it is necessary to follow the path and not be distracted by pseudo-intellectual side trips.


\end{sffamily}\end{footnotesize}
